\section{Proposed Joint Quantization-Aware Training Framework}

This section details the proposed task-oriented communication framework, designed to jointly optimize bit-width allocation at the Access Points (APs) and signal detection at the Central Processing Unit (CPU). As illustrated in Fig. \ref{fig:framework}, the architecture comprises two primary components: a distributed \textit{Resource Inference Module} at each AP and a centralized \textit{Bit-Aware Graph Neural Network (GNN)} at the CPU. The processing pipeline operates as follows: First, each AP performs local linear processing to obtain soft estimates of user symbols. Second, a resource-efficient Policy Network at each AP analyzes the reliability of these estimates to infer a quantization bit-width probability distribution. Third, a Differentiable Quantizer generates a quantized message using a soft relaxation of discrete bit-widths to facilitate gradient backpropagation. Finally, the CPU aggregates these heterogeneous-precision signals via a GNN that weights inputs according to their quantization fidelity.

\begin{figure}[htbp]
\centering
\resizebox{0.95\linewidth}{!}{
\begin{tikzpicture}[
    node distance=1.5cm,
    block/.style={rectangle, draw, fill=blue!10, text width=2.5cm, text centered, rounded corners, minimum height=1.2cm, font=\small},
    sblock/.style={rectangle, draw, fill=green!10, text width=2cm, text centered, rounded corners, minimum height=1cm, font=\small},
    cpublock/.style={rectangle, draw, fill=red!10, text width=3cm, text centered, rounded corners, minimum height=2.5cm, font=\small},
    arrow/.style={-Latex, thick},
    group/.style={rectangle, draw, dashed, inner sep=0.4cm, label=above:\textbf{#1}}
]

% Nodes for AP l
\node (input) {$\mathbf{y}_l, \mathbf{h}_{lk}$};
\node[block, right=0.8cm of input] (lmmse) {Local LMMSE\\Estimation};
\node[sblock, below=0.8cm of lmmse] (policy) {Policy Net $\pi_\phi$\\(Resource Inference)};
\node[block, right=1.5cm of lmmse] (quant) {Differentiable\\Quantizer (LSQ)};

% AP Group Box
\node[group={Access Point $l$}, fit=(input) (lmmse) (policy) (quant)] (ap) {};

% CPU Node
\node[cpublock, right=2.5cm of quant] (gnn) {\textbf{Bit-Aware GNN}\\Attention Aggregation\\+ Transformer};
\node[right=1cm of gnn] (output) {$\hat{\mathbf{s}}$};

% Connections
\draw[arrow] (input) -- (lmmse);
\draw[arrow] (lmmse) -- node[above, font=\footnotesize] {$\hat{s}_{lk}$} (quant);
\draw[arrow] (lmmse.south) -- ++(0,-0.3) -| (policy.north);
\draw[arrow] (policy.east) -| node[right, near start, font=\footnotesize] {$\mathbf{w}_{lk}$} (quant.south);
\draw[arrow] (quant) -- node[above, font=\footnotesize] {$\dot{s}_{lk}$} node[below, font=\footnotesize] {Fronthaul} (gnn);
\draw[arrow] (gnn) -- (output);

\end{tikzpicture}
}
\caption{The proposed Joint Quantization-Aware Training (QAT) framework. Each AP locally infers the optimal bit-width based on instantaneous channel conditions, while the CPU aggregates quantized signals using a bit-aware attention mechanism to compensate for quantization noise.}
\label{fig:framework}
\end{figure}

\subsection{AP-Side Resource Inference Module}
To enable dynamic fronthaul adaptation, each AP $l$ is equipped with a Policy Network $\pi_{\phi}$, parameterized by $\phi$, which determines the optimal quantization precision for each user $k$. Let $\hat{s}_{lk}$ denote the local Linear Minimum Mean Square Error (LMMSE) estimate derived in the system model. The input feature vector to the policy network for the $k$-th user at AP $l$ is defined as $\mathbf{x}_{lk} = [\Re(\hat{s}_{lk}), \Im(\hat{s}_{lk}), \gamma_{lk}]^T \in \mathbb{R}^3$, where $\gamma_{lk} = \|\mathbf{h}_{lk}\|^2 / \sigma^2$ represents the local instantaneous Signal-to-Noise Ratio (SNR). The inclusion of $\gamma_{lk}$ enables the network to differentiate between noise-dominated observations and informative signal components.

The policy network is implemented as a lightweight Multi-Layer Perceptron (MLP) comprising two fully connected layers with ReLU activation. It maps $\mathbf{x}_{lk}$ to a set of unnormalized logits $\mathbf{z}_{lk} \in \mathbb{R}^{|\mathcal{B}|}$, corresponding to the discrete bit-width set $\mathcal{B} = \{0, 2, 4\}$. To facilitate end-to-end training through the discrete selection process, the Gumbel-Softmax relaxation is employed. During the forward pass, a soft selection vector $\mathbf{w}_{lk} \in [0,1]^{|\mathcal{B}|}$ is generated as:
\begin{equation}
    w_{lk, i} = \frac{\exp((\log z_{lk, i} + g_i) / \tau)}{\sum_{j=1}^{|\mathcal{B}|} \exp((\log z_{lk, j} + g_j) / \tau)}, \quad i \in \{1, \dots, |\mathcal{B}|\},
\end{equation}
where $g_i$ are independent samples drawn from the Gumbel$(0, 1)$ distribution, and $\tau > 0$ is the temperature parameter. As $\tau \to 0$, $\mathbf{w}_{lk}$ approaches a one-hot vector representing a discrete choice. This mechanism allows gradients to propagate through the categorical decision of bit-widths during the training phase.

\subsection{Differentiable LSQ Quantizer}
Upon obtaining the selection weights $\mathbf{w}_{lk}$, the signal $\hat{s}_{lk}$ is processed by a bank of quantizers. The Learned Step Size Quantization (LSQ) method is adopted to minimize quantization distortion for non-Gaussian signals. For a target bit-width $b \in \mathcal{B} \setminus \{0\}$, the quantizer $Q_b(\cdot)$ is defined by a learnable step size parameter $\Delta_b$. The quantization of a scalar value $v$ (representing either the real or imaginary part of $\hat{s}_{lk}$) is expressed as:
\begin{equation}
    Q_b(v; \Delta_b) = \Delta_b \cdot \text{round}\left( \text{clip}\left( \frac{v}{\Delta_b}, v_{\min}, v_{\max} \right) \right),
\end{equation}
where $v_{\min} = -2^{b-1}$ and $v_{\max} = 2^{b-1}-1$. To enable gradient descent optimization for the step size $\Delta_b$ and the input $v$, the Straight-Through Estimator (STE) is employed during the backward pass. Specifically, the gradient through the rounding operation is approximated as an identity function within the clipping range.

The final transmitted signal $\dot{s}_{lk}$ is a probability-weighted sum of the outputs from the candidate quantizers:
\begin{equation}
    \dot{s}_{lk} = w_{lk, 0} \cdot 0 + w_{lk, 1} \cdot Q_2(\hat{s}_{lk}) + w_{lk, 2} \cdot Q_4(\hat{s}_{lk}).
\end{equation}
This soft composition allows the system to explore the trade-off between transmitting 0 bits (suppression), 2 bits (coarse quantization), and 4 bits (fine quantization) in a differentiable manner.

\subsection{CPU-Side Bit-Aware GNN Detector}
At the CPU, the aggregated signals exhibit heterogeneous reliability due to the varying quantization precision applied across different APs. To robustly detect user symbols under these conditions, a Bit-Aware GNN is proposed. A fully connected graph is constructed where the $K$ users function as nodes. The feature vector for the $k$-th user node aggregates information from all $L$ APs. Standard aggregation methods (e.g., mean pooling) are suboptimal when signal quality varies significantly.

To address this, a \textit{Bit-Aware Attention} mechanism is introduced. The input feature from AP $l$ for user $k$ is augmented with the expected bit-width $\bar{b}_{lk} = \sum_{j} w_{lk, j} \cdot b_j$. The augmented feature vector is given by $\mathbf{f}_{lk} = [\Re(\dot{s}_{lk}), \Im(\dot{s}_{lk}), \gamma_{lk}, \bar{b}_{lk}]^T$. These features are lifted to a high-dimensional hidden space via an MLP, producing embeddings $\mathbf{h}_{lk}$. The CPU computes attention weights $\alpha_{lk}$ to determine the significance of the contribution from AP $l$:
\begin{equation}
    \alpha_{lk} = \text{Softmax}_l \left( \text{MLP}_{attn}(\mathbf{h}_{lk}) \right).
\end{equation}
The inclusion of $\bar{b}_{lk}$ in the input enables the attention network to assign lower weights to APs that transmitted with 0 or 2 bits and higher weights to those with 4 bits, effectively filtering out quantization noise. The aggregated user feature $\mathbf{u}_k = \sum_{l=1}^{L} \alpha_{lk} \mathbf{h}_{lk}$ is subsequently processed by a Transformer-based encoder to mitigate multi-user interference before the final linear projection to the symbol estimate $\hat{s}_k$.

\subsection{Training Strategy and Complexity}
The framework undergoes end-to-end training. The objective function is a Lagrangian relaxation that balances detection accuracy and fronthaul capacity:
\begin{equation}
    \mathcal{L} = \frac{1}{K} \sum_{k=1}^{K} \| s_k - \hat{s}_k \|^2 + \lambda \left( \frac{1}{LK} \sum_{l,k} \bar{b}_{lk} - C_{\text{avg}} \right)^2,
\end{equation}
where $\lambda$ controls the penalty strength for deviating from the target average capacity $C_{\text{avg}}$. An annealing schedule is employed for the Gumbel temperature $\tau$, decaying from $1.0$ to $0.1$ over the training epochs. This transitions the system from a soft, probabilistic exploration phase to a hard, deterministic selection phase suitable for deployment.

Regarding computational complexity, the AP-side overhead is minimal. The Policy Network requires only $\mathcal{O}(1)$ FLOPs per user-AP pair (specifically, a small MLP with input dimension 3), which is negligible compared to the $\mathcal{O}(N^3)$ complexity of the LMMSE pre-processing. The CPU-side GNN scales linearly with $L$ and quadratically with $K$ due to the attention mechanism, ensuring scalability for dense AP deployments.